\documentclass[11pt]{article}

\usepackage[T1]{fontenc}
\usepackage{lmodern}
\usepackage[margin=1in]{geometry}
\usepackage{amsmath,amssymb,amsthm}
\usepackage{microtype}
\usepackage[hidelinks]{hyperref}
\setlength{\emergencystretch}{2em}

\newtheorem{theorem}{Theorem}[section]
\newtheorem{lemma}[theorem]{Lemma}
\newcommand{\F}{\mathbb F}
\newcommand{\qbinom}[2]{\genfrac{[}{]}{0pt}{}{#1}{#2}_{q}}

\hypersetup{
  pdftitle={On a Claimed Upper Bound for Partial Desarguesian Parallelisms},
  pdfauthor={Carptopus},
  pdfsubject={A uniform correction to an upper bound for partial Desarguesian parallelisms},
  pdfkeywords={Desarguesian spread, partial parallelism, finite geometry, field reduction, Cayley graph}
}

\title{On a Claimed Upper Bound for\\
Partial Desarguesian Parallelisms}
\author{Carptopus\\\small\texttt{carptopus@163.com}}
\date{26 August 2026\\\small v0.1-beta; not peer reviewed}

\begin{document}
\maketitle

\begin{abstract}
Johnson claimed that, when $t$ is prime, a partial Desarguesian
$t$-parallelism in $V(zt,q)$ has at most
\[
\left\lfloor\frac{q^{zt-1}-1}{q^{t-1}-1}\right\rfloor
\]
members, and concluded that no Desarguesian $t$-parallelism exists for odd
prime $t$. We show that this bound fails throughout its nontrivial odd-prime
parameter range. The Cayley-graph construction of Zhang and Zhou uses
$\mathrm{GL}(zt,q)$-images of a standard Desarguesian spread. These images
remain Desarguesian, and their lower bound is strictly larger than Johnson's
claimed upper bound for every prime power $q$, odd prime $t\geq3$, and
$z\geq2$. The comparison is elementary once the two bounds are placed in the
same notation.

We also locate the unsupported step in the proposed correspondence underlying
the bound: affine lines arising from two potentially different extension-field
structures are treated as a common line merely because their defining points
belong to the set-theoretic intersection. Johnson's later book itself records
the associated disjointness issue as unresolved while retaining the claimed
theorem. Earlier work on $PG(5,2)$ already implicitly supplies a complete
finite counterexample, so no priority is claimed for the first counterexample.
A compact certificate of 18 pairwise disjoint Desarguesian $3$-spreads in
$V(6,2)$ is included as an independently checkable finite witness.
\end{abstract}

\noindent\textbf{Keywords.} Desarguesian spread; partial parallelism; finite
geometry; field reduction; Cayley graph.

\medskip
\noindent\textbf{2020 Mathematics Subject Classification.} 51E20, 51A40,
05B25.

\section{Introduction}

Let $V=V(n,q)$ be an $n$-dimensional vector space over $\F_q$, and suppose
that $t\mid n$. A $t$-spread of $V$ is a partition of the nonzero vectors
into $t$-dimensional $\F_q$-subspaces. It is Desarguesian when $V$ can be
regarded as a vector space over $\F_{q^t}$ and the spread members are the
one-dimensional $\F_{q^t}$-subspaces. A partial $t$-parallelism is a family
of $t$-spreads no two of which share a $t$-dimensional subspace.

For $n=zt$, write $D(n,t,q)$ for the largest size of a partial
$t$-parallelism all of whose spreads are Desarguesian. Johnson~\cite[Corollary~1]{JohnsonPaper}
claimed, for prime $t$, the upper bound
\begin{equation}
D(zt,t,q)\leq
J(zt,t,q):=
\left\lfloor
\frac{q^{zt-1}-1}{q^{t-1}-1}
\right\rfloor.
\label{eq:johnson}
\end{equation}
The paper then used~\eqref{eq:johnson} to rule out Desarguesian
$t$-parallelisms for odd prime $t$. An equivalent statement and bound
reappear as Theorem~51 and Corollary~14 of Johnson's
book~\cite[Chapter~10]{JohnsonBook}.

The purpose of this note is not to construct a new large family of spreads.
Instead, it observes that a known construction already contradicts
\eqref{eq:johnson} in every nontrivial odd-prime case. Zhang and
Zhou~\cite{ZhangZhou} derived lower bounds for partial $t$-parallelisms by
taking an independent set in a Cayley graph on $\mathrm{GL}(n,q)$. Their
vertices act on a fixed field-reduction spread. Consequently, every spread in
the resulting family is Desarguesian, even though their extremal notation does
not impose that restriction.

Our main result is the following uniform correction.

\begin{theorem}\label{thm:main}
Let $q$ be a prime power, let $t\geq3$ be an odd prime, let $z\geq2$, and put
$n=zt$. Then
\begin{equation}
D(n,t,q)>
\frac{q^t-1}{q^n-1}\qbinom{n-1}{t-1}
>
\frac{q^{n-1}-1}{q^{t-1}-1}.
\label{eq:main}
\end{equation}
In particular, $D(n,t,q)>J(n,t,q)$, so the upper bound
\eqref{eq:johnson} is false throughout this parameter range.
\end{theorem}

The first inequality in~\eqref{eq:main} comes from the Zhang--Zhou
construction; the second is the new comparison needed to expose the
contradiction.

Section~\ref{sec:orbit} checks that the Zhang--Zhou witnesses satisfy
Johnson's Desarguesian condition. Section~\ref{sec:comparison} proves the
strict comparison. Section~\ref{sec:history} records the earlier finite
counterexample and examines the proof step behind the claimed bound.
Section~\ref{sec:certificate} describes the independent finite certificate.

\section{Desarguesianity of the orbit construction}\label{sec:orbit}

Fix $n$ divisible by $t$. A standard Desarguesian $t$-spread can be obtained
by identifying $V(n,q)$ with $\F_{q^n}$ and taking
\begin{equation}
\mathcal S_0=
\{\,\omega^i\F_{q^t}:
0\leq i<(q^n-1)/(q^t-1)\,\},
\label{eq:standard-spread}
\end{equation}
where $\omega$ is primitive in $\F_{q^n}$. Zhang and
Zhou~\cite[Theorem~3.1]{ZhangZhou} form a Cayley graph on
$\mathrm{GL}(n,q)$ in which two matrices are adjacent precisely when the
corresponding images of $\mathcal S_0$ share a spread member. An independent
set therefore produces a partial $t$-parallelism consisting of spreads
$M\mathcal S_0$.

The following elementary observation places this construction inside the
restricted class counted by $D(n,t,q)$.

\begin{lemma}\label{lem:orbit-desarguesian}
For every $M\in\mathrm{GL}(n,q)$, the spread $M\mathcal S_0$ is
Desarguesian.
\end{lemma}

\begin{proof}
Let
\[
\rho:\F_{q^t}\longrightarrow
\mathrm{End}_{\F_q}(V)
\]
be the extension-field embedding whose one-dimensional orbits give
$\mathcal S_0$. Conjugation defines another embedding
\[
\rho_M(a)=M\rho(a)M^{-1}.
\]
The one-dimensional $\F_{q^t}$-subspaces for the extension structure
$\rho_M$ are exactly the $M$-images of those for $\rho$. Hence
$M\mathcal S_0$ is Desarguesian.
\end{proof}

Johnson's definition requires each spread in the partial parallelism to be
Desarguesian; it does not require all spreads to arise from one common
extension-field embedding. Lemma~\ref{lem:orbit-desarguesian} is therefore
sufficient.

Zhang and Zhou prove~\cite[Corollary~3.4]{ZhangZhou} that their construction
contains more than
\begin{equation}
L(n,t,q):=
\frac{q^t-1}{q^n-1}\qbinom{n-1}{t-1}
\label{eq:L}
\end{equation}
spreads. Since the witnesses used in that proof are the orbit images just
described, Lemma~\ref{lem:orbit-desarguesian} gives
\begin{equation}
D(n,t,q)>L(n,t,q).
\label{eq:DgtL}
\end{equation}

\section{Uniform comparison with the claimed upper bound}
\label{sec:comparison}

Put
\[
J_0(n,t,q)=\frac{q^{n-1}-1}{q^{t-1}-1},
\]
so that $J(n,t,q)=\lfloor J_0(n,t,q)\rfloor$. Expanding the Gaussian
binomial coefficient in~\eqref{eq:L} and cancelling the factor with index
zero gives
\begin{equation}
\frac{L(n,t,q)}{J_0(n,t,q)}
=
\frac{q^t-1}{q^n-1}
\prod_{i=1}^{t-2}
\frac{q^{n-1-i}-1}{q^{t-1-i}-1}.
\label{eq:ratio}
\end{equation}
We prove that this ratio is greater than one.

\begin{lemma}\label{lem:t3}
If $t=3$ and $n=3z$ with $z\geq2$, then $L(n,3,q)>J_0(n,3,q)$.
\end{lemma}

\begin{proof}
Equation~\eqref{eq:ratio} becomes
\[
\frac{L(n,3,q)}{J_0(n,3,q)}
=
\frac{(q^2+q+1)(q^{n-2}-1)}{q^n-1}.
\]
The numerator minus the denominator is
\[
(q^2+q+1)(q^{n-2}-1)-(q^n-1)
=(q+1)(q^{n-2}-q),
\]
which is positive because $n\geq6$.
\end{proof}

\begin{lemma}\label{lem:t4plus}
If $t\geq4$, $n=zt$, and $z\geq2$, then
$L(n,t,q)>J_0(n,t,q)$.
\end{lemma}

\begin{proof}
For $1\leq i\leq t-2$,
\begin{equation}
\frac{q^{n-1-i}-1}{q^{t-1-i}-1}>q^{n-t},
\label{eq:factor}
\end{equation}
because the numerator of the left side exceeds
$q^{n-t}(q^{t-1-i}-1)$. Also,
\begin{equation}
\frac{q^t-1}{q^n-1}
=q^{t-n}\frac{1-q^{-t}}{1-q^{-n}}
>(1-q^{-t})q^{t-n}.
\label{eq:leading}
\end{equation}
Substituting~\eqref{eq:factor} and~\eqref{eq:leading} into
\eqref{eq:ratio} yields
\begin{equation}
\frac{L(n,t,q)}{J_0(n,t,q)}
>(1-q^{-t})q^{(t-3)(n-t)}.
\label{eq:coarse}
\end{equation}
For $t=4$, the right side is minimized at $q=2$ and $n-t=4$, where it is
\[
\frac{15}{16}\,2^4=15>1.
\]
For $t\geq5$, it is minimized at the boundary $q=2$, $t=5$, and $n-t=5$,
where it is
\[
\frac{31}{32}\,2^{10}>1.
\]
Thus the ratio in~\eqref{eq:ratio} is greater than one.
\end{proof}

\begin{proof}[Proof of Theorem~\ref{thm:main}]
By Lemma~\ref{lem:orbit-desarguesian} and the Zhang--Zhou construction,
\eqref{eq:DgtL} holds. Lemma~\ref{lem:t3} handles $t=3$, and
Lemma~\ref{lem:t4plus} handles every odd prime $t\geq5$. Therefore
\[
D(n,t,q)>L(n,t,q)>J_0(n,t,q)
\geq\lfloor J_0(n,t,q)\rfloor,
\]
which is~\eqref{eq:main}.
\end{proof}

The argument actually compares the two numerical bounds whenever $t\geq3$,
$t\mid n$, and $n\geq2t$; primality is retained in
Theorem~\ref{thm:main} because that is the scope of Johnson's claim under
correction.

\section{Historical boundary and the failed proof bridge}
\label{sec:history}

\subsection{An earlier complete finite counterexample}

The parameter triple $(q,t,z)=(2,3,2)$ already shows why the present note
cannot claim the first counterexample. Projective planes in $PG(5,2)$
correspond to three-dimensional vector subspaces of $V(6,2)$. Topalova and
Zhelezova~\cite{Topalova2008} proved that $PG(5,2)$ has only one equivalence
class of projective $2$-spreads under $PGL(6,2)$. A standard field-reduction
spread from $\F_{64}/\F_8$ is Desarguesian. Since
$PGL(6,2)=GL(6,2)$, every $2$-spread in this geometry is therefore
Desarguesian.

Sarmiento~\cite{Sarmiento} had already constructed resolutions of the
$2$-$(63,7,15)$ design associated with $PG(5,2)$; these resolutions are
complete projective $2$-parallelisms. The number of projective planes is
\[
\genfrac{[}{]}{0pt}{}{6}{3}_{2}=1395,
\]
and every spread contains
\[
\frac{2^6-1}{2^3-1}=9
\]
planes. A complete parallelism consequently contains $1395/9=155$ spreads,
all Desarguesian by the uniqueness result. Johnson's claimed bound at these
parameters is only
\[
\left\lfloor\frac{2^5-1}{2^2-1}\right\rfloor=10.
\]
Thus the older literature already implicitly contains the contradiction
$155>10$. Later work of Topalova and Zhelezova~\cite{Topalova2010} gives
further explicit constructions and classifications in the same geometry.

\subsection{Where the proposed correspondence loses a hypothesis}

Johnson derives~\eqref{eq:johnson} from a correspondence between partial
Desarguesian $t$-parallelisms and translation nets assembled from rational
Desarguesian partial spreads~\cite[Theorem~3]{JohnsonPaper}. In the proof, two
affine planes $\pi_1,\pi_2$ arise from two Desarguesian spreads. These spreads
may use different copies of $\F_{q^t}$ inside
$\mathrm{End}_{\F_q}(V)$.

For two points $P,Q\in\pi_1\cap\pi_2$, the proof then treats the affine line
through $P$ and $Q$ as a line common to both planes. Uniqueness of a line
through two points holds within each affine-plane incidence structure, but it
does not identify the two lines when the extension-field structures differ.
The same issue occurs when $0$ and $Q-P$ are used to infer a common component.
Consequently, the set-theoretic intersection $\pi_1\cap\pi_2$ has not been
shown to be an affine subplane of either plane, and the subsequent subfield
dimension argument lacks its required premise.

The book version makes the status of this bridge especially clear. Before
restating the result, Johnson~\cite[pp.~134--136]{JohnsonBook} says that it is
``not clear'' whether the two rational Pappian partial spreads are disjoint
outside their common regulus. Theorem~51 nevertheless restates the
equivalence, leaves the common-affine-subplane assertion to the reader, and
Corollary~14 repeats the upper bound.

The numerical contradiction in Theorem~\ref{thm:main} proves that this
unresolved bridge cannot hold under the stated hypotheses in general. This
note does not attempt to classify additional compatibility assumptions under
which a modified correspondence might be valid.

\section{A compact independent certificate}\label{sec:certificate}

The accompanying verification data give 18 explicit matrices in
$\mathrm{GL}(6,2)$. Starting from the field-reduction spread
$\F_{64}/\F_8$, the checker reconstructs the 18 image spreads and verifies
that
\begin{enumerate}
\item every matrix is invertible;
\item every image is a $3$-spread of $V(6,2)$;
\item the 18 spreads are distinct; and
\item every pair is disjoint as a set of three-dimensional subspaces.
\end{enumerate}

Every image is Desarguesian by Lemma~\ref{lem:orbit-desarguesian}, so the
certificate independently establishes
\[
D(6,3,2)\geq18>10.
\]
It is intentionally smaller than the old complete $155$-spread witness. Its
purpose is reproducibility: checking the correction does not require
reconstructing the earlier full classifications or parallelisms. The checker
also contains stable destructive controls which reject a singular map and an
intentionally duplicated spread.

From this public entry directory, run
\begin{verbatim}
python -X utf8 verification\check_desarguesian_certificate.py
\end{verbatim}

The finite certificate is supplementary evidence. The general result in
Theorem~\ref{thm:main} rests on the symbolic comparison in
Section~\ref{sec:comparison}, not on finite enumeration.

\section{Scope of the correction}

The mathematical contribution of this note is deliberately limited:
\begin{itemize}
\item it identifies the Zhang--Zhou orbit lower bound as a uniform
contradiction to Johnson's claimed upper bound over the whole nontrivial
odd-prime range;
\item it explains why all orbit witnesses meet Johnson's individual
Desarguesianity requirement;
\item it shows that the unresolved common-subplane bridge is fatal to the
claimed theorem under its published hypotheses; and
\item it supplies a compact independent finite witness.
\end{itemize}

The note does not claim the first finite counterexample, the first
Desarguesian $3$-parallelism, the first observation that the common-subplane
step needs proof, or a best corrected replacement for Johnson's
correspondence. It also does not claim external expert confirmation or peer
review.

\section*{Acknowledgement of AI assistance}

The author used OpenAI Codex for literature triage, symbolic checks, proof
stress testing, verification-code review, and drafting assistance. The author
reviewed the mathematical argument and assumes responsibility for the claims
and presentation.

\begin{thebibliography}{9}

\bibitem{JohnsonPaper}
N.~L. Johnson,
\emph{The non-existence of desarguesian t-parallelisms, t an odd prime},
Note di Matematica 30 (2010), no.~1, 13--16.
\url{https://doi.org/10.1285/i15900932v30n1p13}

\bibitem{JohnsonBook}
N.~L. Johnson,
\emph{Combinatorics of Spreads and Parallelisms},
CRC Press, 2010, Chapter~10, Theorem~51 and Corollary~14.
\url{https://doi.org/10.1201/EBK1439819463}

\bibitem{Sarmiento}
J.~F. Sarmiento,
\emph{On point-cyclic resolutions of the 2-(63,7,15) design associated with
PG(5,2)},
Graphs and Combinatorics 18 (2002), 621--632.
\url{https://doi.org/10.1007/s003730200046}

\bibitem{Topalova2008}
S.~T. Topalova and S.~D. Zhelezova,
\emph{Classification of the 2-spreads of PG(5,2)},
Proceedings of the Thirty Seventh Spring Conference of the Union of Bulgarian
Mathematicians (2008), 163--168.
\url{https://www.math.bas.bg/smb/2008_PK/2008/pdf/163-168.pdf}

\bibitem{Topalova2010}
S.~Topalova and S.~Zhelezova,
\emph{2-spreads and transitive and orthogonal 2-parallelisms of PG(5,2)},
Graphs and Combinatorics 26 (2010), 727--735.
\url{https://doi.org/10.1007/s00373-010-0943-8}

\bibitem{ZhangZhou}
T.~Zhang and Y.~Zhou,
\emph{New lower bounds for partial k-parallelisms},
Journal of Combinatorial Designs 28 (2020), 75--84.
\url{https://doi.org/10.1002/jcd.21681}

\end{thebibliography}

\end{document}
